% Source: https://artofproblemsolving.com/community/c3h1597305_interesting_proof_problem
% Stable identity: aops-topic-1597305-q01
\textbf{(a)} Since $a^2+b^2$ is even, $a^2$ and $b^2$ have the same parity, so $a$ and $b$ have the same parity. Hence $a+b$ and $a-b$ are both even, and
\[
c=\dfrac{a+b}{2},\qquad d=\dfrac{a-b}{2}
\]
are integers. Now
\[
c^2+d^2=\dfrac{(a+b)^2+(a-b)^2}{4}=\dfrac{2a^2+2b^2}{4}=\dfrac{a^2+b^2}{2}=\dfrac {x}{2} .
\]
So $\dfrac {x}{2}=c^2+d^2$. $\blacksquare$

\textbf{(b)} Take $x=2=1^2+1^2$. Its half is $1$, and $1$ is not a sum of two squares of positive integers (such a sum is at least $1+1=2$). So with $0$ excluded the statement fails; the formula above gives $1=1^2+0^2$.

If $a,b>0$ and $a\ne b$, we may assume $a>b$. Then $c=\dfrac{a+b}{2}>0$ and $d=\dfrac{a-b}{2}>0$, so $\dfrac {x}{2}=c^2+d^2$ is a sum of two squares of positive integers. (For example, $x=10=3^2+1^2$ gives $5=2^2+1^2$.) $\blacksquare$
